\documentclass[11pt]{article}
\usepackage{amsmath,amssymb,amsthm}
\usepackage[margin=1in]{geometry}
\newtheorem{theorem}{Theorem}
\newtheorem{lemma}[theorem]{Lemma}
\newtheorem{proposition}[theorem]{Proposition}
\newtheorem{conjecture}[theorem]{Conjecture}
\theoremstyle{remark}
\newtheorem{remark}[theorem]{Remark}
\newcommand{\Q}{\mathbb Q}
\newcommand{\Z}{\mathbb Z}
\newcommand{\C}{\mathbb C}
\title{P5: is the growth series $F_T$ of $W_T$ $D$-finite? A reduction, not a resolution}
\date{2026-09-26}
\begin{document}
\maketitle

\noindent Labels: PROVED means a complete argument, given the inputs it cites. VERIFIED means
numerical/computational evidence with the precision or range stated (no error control unless said
otherwise). HEURISTIC means an argument without a closed error bound, offered as an informed guess.
Notation follows \texttt{paper/journal/paper1.tex} (``paper1'') for $W_T$ and
\texttt{paper/journal/paper2a.tex} (``paper2a'') for $U,V$ and the travel operator.

\section*{Summary of outcome}

This note does \emph{not} resolve Conjecture~\texttt{conj:WT-nonDfinite} of paper1 (that $F_T$, the
growth series of the rational right-triangle group $W_T$, is not $D$-finite). What it does:
\begin{itemize}
\item States precisely what is already PROVED (Theorem~\texttt{thm:WT-lowcomplexity}: no relation of
complexity $\le d$ in explicit finite boxes) and what is only VERIFIED/HEURISTIC (the linear carry-width
growth, the non-transfer of the travel-pole mechanism), reading paper1 \S8.7--8.9 in full.
\item Isolates the exact mechanism that proves non-$D$-finiteness for $U,V$ in paper2a
(Theorem~\texttt{thm:nonDfinite}, via Lemma~\texttt{lem:odeposes}) and states, as a genuinely new but
easy PROVED lemma (\S\ref{sec:reduction}), the precise sufficient condition under which the same
mechanism would give non-$D$-finiteness of $F_T$: a meromorphic continuation of $F_T$ to a domain
containing infinitely many of its poles. This is an immediate corollary of the already-proved
domain-general Lemma~\texttt{lem:odeposes}; the content is the reduction, not new analytic input.
\item Explains precisely why that hypothesis is not currently available for $F_T$ (\S\ref{sec:gap}):
unlike $U,V,G^T_0$, $F_T$ has no known continuation past its radius of convergence at all, so the
sufficient condition of \S\ref{sec:reduction} cannot even be tested for $F_T$ with tools now in the
repository.
\item Records the FKS/Akiyama--Frougny analogy in the precise form paper1 already uses it
(\S\ref{sec:analogy}), and states as HEURISTIC (paper1's own word) why the rational-base transfer
picture is expected to fail for $W_T$: the numeration base $\zeta_T$ has $|\zeta_T|=1$, so nothing
contracts the carry, and the carry-width data computed by \texttt{wt\_carry} shows linear growth in the
sampled range, not the boundedness that drives the $\mathrm{BS}(p,q)$, $p\mid q$ rationality proof or
the Akiyama--Frougny--Sakarovitch automaticity results.
\item States the remaining gap as three independently missing ingredients (\S\ref{sec:gap}), none of
which is closed here.
\end{itemize}
The row stays OPEN. No new computation was run for this note beyond re-deriving already-published
numbers from the existing \texttt{wt\_carry}/\texttt{wt\_growth} output files, which were re-read but
not regenerated (see \S\ref{sec:verif}).

\section{What $W_T$ and $F_T$ are}\label{sec:def}

Fix a right triangle $T$ with positive, unequal rational legs, encoded (paper1 \S8) by coprime
integers $(c,e)$ with $0<e<2c$ via the minimal polynomial $\mu_T=ct^2-et+c$ of the rotation number
$\zeta_T$ (a primitive root of $\mu_T$, $|\zeta_T|=1$, $\zeta_T\notin\R$). The group
\[
  W_T=\langle R_0,R_1,R_2\rangle\le\mathrm{Isom}(\R^2)
\]
is the image of the universal reflection group $W_{\mathrm{gen}}$ under the representation
$\rho_T$ specializing the generic rotation parameter to $\zeta_T$ (paper1, \S2283--2340: $W_T$ is
dense in $\mathrm{Isom}(\R^2)$, linear over $\Q$, solvable of derived length exactly $3$, not finitely
presented, quasi-isometric to $\R^2\times$ a horocyclic product of Baumslag--Solitar-type spaces). Its
\emph{growth series} with respect to the standard generating set $\{R_0,R_1,R_2\}$ is
\[
  F_T(x)=\sum_{d\ge0}u_d^T\,x^d,\qquad u_d^T=\#\{g\in W_T:\ \ell_T(g)=d\},
\]
paper1 line 2548. By Theorem~\texttt{thm:universality-sharp} the layers $u_d^T$ agree with the
$T$-independent \emph{universal} layers $u_d$ (of $W_{\mathrm{gen}}$/its associated
$\mathrm{Elt}$-model) up to a shape-dependent onset depth $n_T\ge c_T$
(Corollary~\texttt{cor:kernel-classification-restate}), and only deviate from them beyond that depth.

\paragraph{The $D$-finiteness question (P5).} $F_T$ is $D$-finite (holonomic) if it satisfies a linear
ODE $\sum_{i=0}^r Q_i(x)F_T^{(i)}(x)=0$ with polynomial coefficients $Q_i$, $Q_r\not\equiv0$; equivalently
(Stanley, \emph{Enumerative Combinatorics} II, \S6.4 — cited as in paper1/paper2a, not independently
re-verified in this note) $(u_d^T)$ is P-recursive. Conjecture~\texttt{conj:WT-nonDfinite} (paper1, line
2734) asserts $F_T$ is \emph{not} $D$-finite, for every such $T$, and likewise for the ``toy groups''
$\Z[\zeta_T^{\pm1}]\rtimes_{\zeta_T}\Z$ with generators $\{1,\mathrm{shift}\}$.

\section{What is already proved: a computer-assisted finite-order exclusion}\label{sec:proved}

\begin{itemize}
\item[(PROVED, computer-assisted)] Theorem~\texttt{thm:WT-lowcomplexity} (paper1, line 2564): for every
such $T$, $F_T$ satisfies no algebraic relation $\sum_iP_iF_T^i=0$ with $(k,\max\deg P_i)$ at most
$(1,14),(2,9),(3,6),(4,5),(5,4),(6,3)$, and $(u_d^T)$ satisfies no P-recursion with $(r,\max\deg Q_i)$
at most $(1,14),\dots,(11,0)$ (paper1 line 2572--2574). The proof uses only the $33$ shape-independent
universal terms $u_0,\dots,u_{32}$ (certified through $u_{22}$ by a \texttt{native\_decide} Lean
computation, tier (ii) trust; through $u_{32}$ by an exact Rust BFS plus modular-plus-exact certificate,
tier (iii) trust — see paper1 \S8.6/App.~\texttt{sec:reproduce}) and full-rank-mod-$(2^{61}-1)$ linear
algebra in \texttt{code/zeta\_probe/wt\_growth/guess.py}, which is exact in the direction used (full
rank mod $p$ implies full rank over $\Q$) but is a script, not formalized.
\item[(PROVED, conditional)] Theorem~\texttt{thm:WT-blowup} (paper1, line 2625): \emph{if} the universal
series $U$ is not $D$-finite (paper2a's Theorem~\texttt{thm:nonDfinite}, unconditional there — see
\S\ref{sec:mechanism}), then for every complexity bound $d$ only finitely many similarity classes of
$T$ have $F_T$ of complexity $\le d$. The proof is a clean finite-dimensional nested-kernel argument and
is ineffective (ends at line 2670: ``The proof is ineffective: it gives no value of $N(d)$'').
\item[(explicit calibration against a false positive, PROVED)] Remark~\texttt{rem:BS-calibration}
(paper1, line 2604): the structurally similar Baumslag--Solitar group $\mathrm{BS}(1,2)$ \emph{does} have
a rational growth series (Collins--Edjvet--Gill, cited), of type $(14,13)$, invisible to a
finite-horizon search until $\ge29$ terms are available. This is why
Theorem~\texttt{thm:WT-lowcomplexity}, a finite-order exclusion, cannot by itself be read as evidence
for the conjecture: a rational $F_T$ of moderate type would currently be undetected.
\end{itemize}
None of the above proves Conjecture~\texttt{conj:WT-nonDfinite}. It proves an exclusion of relations in
explicit finite boxes, and a conditional statement whose hypothesis is exactly (an instance of) the
conjecture's cousin for the universal series $U$ — itself proved unconditionally in paper2a, but by a
mechanism (\S\ref{sec:mechanism}) that paper1 explicitly records as not transferring to $F_T$
(paper1, line 2766: ``does not appear to transfer'').

\section{The mechanism that works for $U,V,G^T_0$}\label{sec:mechanism}

Theorem~\texttt{thm:nonDfinite} of paper2a proves $U,V$, and the travel resolvent $G^T_0$ are not
$D$-finite, via:
\begin{enumerate}
\item[(a)] a meromorphic continuation of the series to the open unit disc $D=\{|x|<1\}$ (for $G^T_0$,
unconditionally, from the Fredholm-determinant representation of the travel operator
$\mathcal T=D_eK$, $K=[q^{\max(s,s')}]_{s,s'\ge0}$; for $U,V$, standing on Theorems~\texttt{thm:model}
and~\texttt{thm:RJ});
\item[(b)] infinitely many poles of that continuation inside $D$ (for $G^T_0$ unconditional — every
$q_m$ in $q$, every $\pm x_m$ in $x$; for $U,V$, at every $x_m$ resp.\ $\pm x_m$, from
Theorems~\texttt{thm:V},~\texttt{thm:U} and Propositions~\texttt{prop:Vminusx},~\texttt{prop:minusx});
\item[(c)] Lemma~\texttt{lem:odeposes} (paper2a, line 1468), a fully general and elementary fact: if $f$
is meromorphic on a domain $\Omega$ and solves $\sum_{j\le r}a_jf^{(j)}=b$ off its poles, with
$a_0,\dots,a_r,b$ holomorphic on $\Omega$ and $a_r\not\equiv0$, then every pole of $f$ in $\Omega$ is a
zero of $a_r$. Since $a_r$ is holomorphic and not identically zero, its zero set in any relatively
compact subdomain is finite. Infinitely many poles in $D$ (from (a),(b)) then contradicts a would-be
holomorphic-coefficient ODE on a neighbourhood of $\overline D$, hence the coefficients cannot be
polynomial either (a polynomial is holomorphic everywhere).
\end{enumerate}
This is the only mechanism the repository currently has for proving non-$D$-finiteness of a series of
this general shape. Any attempt to prove Conjecture~\texttt{conj:WT-nonDfinite} by the same route needs
(a) and (b) to hold for $F_T$ itself (or for some series determining $F_T$, e.g.\ the toy groups' growth
series).

\section{A trivial but honest reduction}\label{sec:reduction}

\begin{lemma}[Reduction, PROVED — immediate from Lemma~\texttt{lem:odeposes}]\label{lem:reduction}
Fix $T$. Suppose $F_T$ (equivalently, its analytic continuation from a neighbourhood of $0$, if one
exists) extends to a meromorphic function on some domain $\Omega\subseteq\C$ containing $0$, and that
this continuation has infinitely many poles lying in some relatively compact subset of $\Omega$. Then
$F_T$ is not $D$-finite: it satisfies no linear ODE with coefficients holomorphic on $\Omega$
(a fortiori none with polynomial coefficients, since a nonzero polynomial has finitely many zeros in
$\C$, hence in any relatively compact set, unless it has none there at all — the same finiteness
argument applies verbatim).
\end{lemma}
\begin{proof}
Identical to the proof of Theorem~\texttt{thm:nonDfinite} in paper2a: apply Lemma~\texttt{lem:odeposes}
with $\Omega$ in place of the unit disc. A pole of $F_T$ in $\Omega$ that is not a zero of the leading
coefficient $a_r$ is impossible, and $a_r$ has only finitely many zeros in any relatively compact subset
of $\Omega$ (identity theorem, $a_r\not\equiv0$ holomorphic), contradicting the assumed infinitude of
poles there.
\end{proof}

\noindent This lemma has zero new mathematical content beyond paper2a's own argument; it is stated here
only to make explicit that \emph{the entire missing step of P5 is producing hypotheses (a) and (b) for
some domain $\Omega$} — nothing else is needed once they are available. This sharpens ``prove
Conjecture~\texttt{conj:WT-nonDfinite}'' to ``find any domain and any continuation of $F_T$ into it with
infinitely many poles there,'' which is the same reduction paper2a already performs for $U,V,G^T_0$, now
made explicit as the target for $F_T$.

\section{Why the hypotheses of Lemma~\ref{lem:reduction} are not available}\label{sec:gap}

Paper1 (\S8.8, lines 2732--2774) already gives the precise reason the $U,V,G^T_0$ route does not
obviously transfer, which this note restates without weakening it:

\begin{enumerate}
\item[\textbf{Gap 1 (no known continuation).}] Unlike $U,V,G^T_0$, there is no known representation of
$F_T$ analogous to the Fredholm determinant of the travel operator (item (a) of \S\ref{sec:mechanism}).
$F_T$ is presently known only as a power series from finite BFS layers (\texttt{wt\_growth}) plus the
universality identification with $U$'s layers up to depth $n_T$. Nothing in the repository establishes
that $F_T$ continues past its radius of convergence $1/\limsup(u_d^T)^{1/d}$ at all. Hypothesis (a) of
Lemma~\ref{lem:reduction} is therefore currently untestable for $F_T$, not merely unproved.
\item[\textbf{Gap 2 (the natural transfer variable is different, and its dynamics look wrong).}] The
combinatorial source of $U,V$'s infinitely many poles is the travel block, a sub-transfer with
catalytic variable that is \emph{not} damped: paper1 identifies the corresponding structure for $W_T$ as
the \emph{carry} $f\in\Z[t^{\pm1}]$ in the identification $A\sim A+2\mu_Tf$ of edge profiles giving the
same group element (Corollary~\texttt{cor:kernel-classification-restate}). Long division of a profile
by $2\mu_T$ from the low end realizes numeration in base $\zeta_T$, $|\zeta_T|=1$. Because the base does
not contract ($|\zeta_T|=1$, as opposed to a rational base $q/p>1$ that contracts a remainder to a
bounded window), nothing forces the carry state to stay in a bounded set: it performs a walk in the
plane $\Z[\zeta_T]$, not a bounded automaton state.
\item[\textbf{Gap 3 (a proof would need a carry-catalytic transfer equation, and none exists).}]
Paper1 states, explicitly labelled HEURISTIC there (line 2771, ``We have not proved this, and state it
only as a heuristic''): the travel-pole mechanism of $U,V$ is believed to be collapsed by the carry
mechanism, i.e.\ pure-travel elements remain distinct in $W_T$ but the \emph{magnitude sector} of the
transfer producing the travel poles does not survive intact once carries are accounted for. No
functional/transfer equation for $F_T$ with the planar carry as an explicit catalytic variable is
known to exist (paper1, line 2772--2773: ``A proof of Conjecture~\texttt{conj:WT-nonDfinite} would seem
to need a transfer in which the planar carry is the catalytic variable, and no such equation is
available.''). This is the mathematical crux: without such an equation there is no candidate route to
either hypothesis (a) or (b) of Lemma~\ref{lem:reduction} for $F_T$.
\end{enumerate}

\noindent Closing P5 unconditionally therefore requires, at minimum, either (i) an explicit meromorphic
continuation of $F_T$ together with a proof of infinitely many poles inside its domain (the direct
route, blocked by Gaps 1--3), or (ii) some entirely different obstruction to $D$-finiteness not based on
pole confinement — e.g.\ a genuinely unbounded-state automaticity/numeration argument transferred from
Akiyama--Frougny--Sakarovitch (\S\ref{sec:analogy}), which would need to be developed from scratch, since
their theorems are stated for contracting (rational, $>1$) bases and $|\zeta_T|=1$ is exactly the
boundary case they do not cover.

\section{The FKS / Akiyama--Frougny analogy, read precisely}\label{sec:analogy}

Paper1 (line 2741--2750) draws the analogy this note is asked to examine. Restated precisely, with the
literature flagged (\S\ref{sec:lit}):
\begin{itemize}
\item Freden--Knudson--Schofield (cited in paper1's bibliography as
\texttt{bibitem\{FKS2011\}}, ``Growth in Baumslag--Solitar groups'' or similarly titled — \textbf{FKS
here stands for the surname initials of these three authors, Freden, Knudson, Schofield; it is not an
acronym for Furstenberg or Flajolet--Kurshan--Steele}, contrary to one plausible guess raised in the
task prompt) computed the growth of the horocyclic subgroup of $\mathrm{BS}(p,q)$ and proved it rational
when $p\mid q$.
\item For $p\nmid q$ (first case $\mathrm{BS}(2,3)$): no geodesic combing of that subgroup is
context-free, and the growth series is conjectured (not proved, per paper1) irrational; the underlying
arithmetic is numeration in the rational base $q/p$, whose integer expansions form a non-regular
language, citing Akiyama--Frougny--Sakarovitch (\texttt{bibitem\{AFS2008\}} in paper1, ``Powers of
rationals modulo 1 and rational base number systems'' or similarly titled).
\item Paper1 explicitly calls Conjecture~\texttt{conj:WT-nonDfinite} ``the unit-circle counterpart of
that situation, and \dots\ at least as hard'' (line 2749--2750). The analogy is therefore stated in
paper1 itself as a difficulty comparison, not as a proof strategy with a known target theorem to invoke:
the $p\nmid q$ non-rationality of $\mathrm{BS}(p,q)$'s horocyclic growth series is \emph{itself an open
problem} in the cited literature (per paper1's own account), so P5's ``next step'' points at a technique
whose home instance is unsolved.
\end{itemize}
\textbf{Consequence for P5.} The suggested approach (``bounded-state structure / rational-base
numeration'') is exactly the numeration-in-base-$\zeta_T$ picture of Gap 2 above, and paper1 already
tried it as far as the repository's tools allow: it looked for a bounded carry window (which would be
the $W_T$ analogue of a contracting rational base giving a finite automaton) and found none, in the
numerically sampled range (below).

\section{The carry-width numerics, re-read (not regenerated)}\label{sec:verif}

\begin{itemize}
\item[(VERIFIED, in the narrow sense of ``re-read the existing output,'' not re-run)]
\texttt{code/zeta\_probe/wt\_growth/src/bin/wt\_carry.rs} computes, for $(c,e)=(2,1)$ and $d\le32$, the
exact optimal carry width (min-plus DP over carries $f$, exact \texttt{i128} arithmetic in
$\Z[1/c][\zeta]$) needed to realize the $W_T$-geodesic length of every element of the universal ball of
radius $d$. Paper1 (line 2761--2763) reports this width as $\lfloor(d-12)/2\rfloor$ for all $d\le32$ —
linear growth, with no plateau, in the sampled range. This is exactly the opposite of the boundedness a
finite-state/bounded-carry automaton picture (the P5 ``next step'') would need, in the one shape and
depth range where it has been checked. Paper1 itself labels this ``verified, not proved'' (line
2764--2765 and again at line 23 of \texttt{wt\_growth/README.md}: ``The carry-window law is verified,
not proved.''). I did not re-run \texttt{wt\_carry} for this note (no new computation was judged
necessary or worth the disk/CPU budget for a result already on file and already correctly labelled); I
only re-read \texttt{code/zeta\_probe/wt\_growth/README.md} and the cited line range of paper1.tex to
confirm the claim and its status label match.
\item No sanity Rust/Arb script is added by this note. A genuinely new contribution here would be
running \texttt{wt\_carry} at additional $(c,e)$ shapes and larger $d$ to test whether the linear law
$\lfloor(d-12)/2\rfloor$ is shape-independent or an artifact of $(c,e)=(2,1)$; this was judged out of
scope for this reduction note (it would extend, not close, the existing VERIFIED claim) and is recorded
as the first concrete next action in \S\ref{sec:next}.
\end{itemize}

\section{Open gap / weakest points, stated explicitly}\label{sec:weak}

\begin{enumerate}
\item The single mathematical gap that would close P5 by the pole-confinement route is exactly Gap 1
of \S\ref{sec:gap}: \emph{no domain $\Omega$ and no meromorphic continuation of $F_T$ into $\Omega$ is
known}. Lemma~\ref{lem:reduction} (this note) shows this is the \emph{only} gap for that route — once
(a) and (b) are supplied, non-$D$-finiteness follows by the identical (already-proved) argument as for
$U,V,G^T_0$.
\item The weakest point of the existing repository evidence is Theorem~\texttt{thm:WT-lowcomplexity}
being read as more than it is: it is a finite-order exclusion, explicitly calibrated against a false
positive ($\mathrm{BS}(1,2)$, Remark~\texttt{rem:BS-calibration}) that a search of comparable size would
have missed. It gives no asymptotic information and cannot, by itself, ever prove
Conjecture~\texttt{conj:WT-nonDfinite} (a conjecture about all orders, not any finite box).
\item The carry-width linear growth (\S\ref{sec:verif}) is VERIFIED at one shape, one depth range
($d\le32$, $(c,e)=(2,1)$), and paper1 itself does not claim it generalizes; it is not even a HEURISTIC
argument for non-$D$-finiteness by itself, only evidence against the one concrete route (bounded carry
automaton) that P5's own suggested next step names.
\item The FKS/Akiyama--Frougny analogy (\S\ref{sec:analogy}) is, on inspection, an analogy to an
\emph{open} problem ($\mathrm{BS}(p,q)$, $p\nmid q$, non-rationality), not a solved one whose method could
be imported. This considerably weakens ``Bounded-state structure / Rational-base numeration'' as a
stated next step: even a full transfer of the Akiyama--Frougny apparatus to $W_T$ would, at best, reduce
P5 to the currently-open $\mathrm{BS}(p,q)$ problem's own difficulty, not solve it.
\item This note performed no new computation and no new proof beyond Lemma~\ref{lem:reduction} (itself
an immediate corollary of an existing lemma). It should not be read as progress on the conjecture, only
as a precise map of where the difficulty is and an explicit, minimal target (Gap 1) for future work.
\end{enumerate}

\section{Concrete next actions (not carried out here)}\label{sec:next}
\begin{enumerate}
\item Run \texttt{wt\_carry} at additional shapes ($(c,e)=(5,-6)$ i.e.\ $(1,2)$, $(53,-90)$ i.e.\ $(2,7)$,
and one or two more laboratory shapes) and larger $d$ (subject to the RSS/wall caps in
\texttt{wt\_growth/lab.sh}, \texttt{lab2.sh}, i.e.\ $\le3000$MB, $\le285$s per run under
\texttt{runcap.sh}) to see whether the carry width's linear law $\lfloor(d-12)/2\rfloor$ is
shape-universal or shape-dependent. A shape-dependent law would be interesting; a universal one would
still not prove non-$D$-finiteness (Gap 3), but would strengthen the heuristic.
\item Attempt to write \emph{any} functional equation for $F_T$ or for the toy-group growth series
$\sum u_d^{\mathrm{toy}}x^d$ with the planar carry as an explicit catalytic (extra) variable, even a
conjectural one to be verified only numerically at first — this is the concrete missing object named at
Gap 3, and the toy groups (already computed exactly, not modularly, in
\texttt{code/zeta\_probe/zeta\_toy\_growth/}) are the more tractable target since their growth series are
known exactly rather than only through a universality bound.
\item If a transfer equation is found, check whether it exhibits a Fredholm-determinant structure whose
zero set can be analyzed the way \texttt{thm:travelres} analyses $G^T_0$; if so, Lemma~\ref{lem:reduction}
of this note applies immediately and would close P5.
\end{enumerate}

\section{Literature relied on — memory flags}\label{sec:lit}

\begin{itemize}
\item \textbf{Independently checked in this session}: every paper1.tex and paper2a.tex line/label cited
above was read directly from the files in this repository (\texttt{paper/journal/paper1.tex},
\texttt{paper/journal/paper2a.tex}) during this session; the \texttt{wt\_growth} README and
\texttt{wt\_carry.rs} source comments were likewise read directly.
\item \textbf{Cited from paper1's own bibliography, not independently verified against the primary
source in this session} (i.e.\ I read paper1's citation of them, but did not fetch or check
Freden--Knudson--Schofield or Akiyama--Frougny--Sakarovitch's papers themselves): the precise theorem
statements attributed to \texttt{FKS2011} and \texttt{AFS2008} above are paper1's paraphrase, taken at
face value. The identification of ``FKS'' with Freden--Knudson--Schofield rather than some other
expansion is read off paper1's \texttt{\textbackslash bibitem\{FKS2011\}} entry
(\texttt{paper/journal/paper1.tex}, line $\sim$5352: ``E.~M.~Freden, T.~Knudson and J.~Schofield,
\emph{Growth in \dots}''), which \emph{was} read directly in this session, so this identification itself
is independently checked, not a memory guess.
\item \textbf{From the assistant's training-data memory, not checked against any source in this
session, and flagged as such}: (i) the general fact that Baumslag--Solitar groups $\mathrm{BS}(p,q)$
with $p\mid q$ have rational growth series and $p\nmid q$ is open/conjectured non-rational — this
matches paper1's own statement and is not contradicted by it, but I did not independently verify it
against Collins--Edjvet--Gill or Freden--Knudson--Schofield directly; (ii) any general background claim
about the Akiyama--Frougny--Sakarovitch rational-base numeration system being a non-regular language and
the associated automaticity theory — I recall this area exists and roughly what it says, but did not
re-derive or re-check any specific theorem statement from it beyond what paper1 already states; this
note's Gap 2/Gap 3 analysis does not depend on the precise form of any AFS theorem, only on the
structural observation (independently verifiable from the $(c,e)=(2,1)$ carry data above) that
$|\zeta_T|=1$ gives a non-contracting base, which is paper1's own point, not a citation-dependent one.
\item \textbf{Stanley, \emph{Enumerative Combinatorics} vol.\ II, \S6.4} (equivalence of $D$-finite
series and P-recursive coefficients): cited exactly as paper1 and paper2a cite it; not independently
re-derived or checked against the book in this session.
\end{itemize}

\end{document}
